\documentclass[11pt]{article}
\usepackage[T1]{fontenc}
\usepackage{amsmath,amssymb,amsthm,microtype,booktabs}
\usepackage[margin=1in]{geometry}

\newtheorem*{theorem}{Theorem}
\newtheorem*{lemma}{Lemma}
\newcommand{\N}{\mathbb N}
\newcommand{\vthree}{\nu_3}

\begin{document}

\begin{center}
{\Large\bfseries A four-dimensional unit-step walk with no four collinear vertices}
\end{center}

\begin{theorem}
There is an infinite walk in $\N^4$, using only the four standard basis
vectors as steps, with no four collinear vertices. Equivalently, there is
an infinite four-letter word with no weak abelian cube.
\end{theorem}

\begin{proof}
Write the ternary expansion of $n$ as $n=\sum_{j\ge0}d_j3^j$, with
$d_j\in\{0,1,2\}$, and put
\[
 a_n=(-1)^{d_1+d_3+d_5+\cdots},\qquad
 b_n=(-1)^{d_0+d_2+d_4+\cdots}.
\]
Let
\[
 X_n=\sum_{j<n}a_j,\qquad Y_n=\sum_{j<n}b_j,
 \qquad R_n=(X_n,Y_n,n)\in\mathbb Z^3.
\]
Appending a ternary digit gives, for $r=0,1,2$,
\begin{equation}\label{eq:rec}
 a_{3n+r}=b_n,\qquad b_{3n+r}=(-1)^r a_n,
\end{equation}
and hence
\begin{equation}\label{eq:prefix}
 X_{3n+r}=3Y_n+r b_n,\qquad
 Y_{3n+r}=X_n+\varepsilon_r a_n,
 \quad (\varepsilon_0,\varepsilon_1,\varepsilon_2)=(0,1,0).
\end{equation}

We first record the only arithmetic fact we need.

\medskip
\noindent\textbf{Lemma.}
If $m<n$ and $(a_m,b_m)=(a_n,b_n)$, then
\begin{equation}\label{eq:valuation}
 \vthree\!\left((X_n-X_m)^2-3(Y_n-Y_m)^2\right)=\vthree(n-m),
\end{equation}
and the quadratic expression is nonzero.

\smallskip
Indeed, write $m=3u+r$ and $n=3v+s$. If $r\ne s$, equality of the
first signs in \eqref{eq:rec} gives $b_u=b_v$, so by \eqref{eq:prefix}
\[
 X_n-X_m\equiv(s-r)b_u\not\equiv0\pmod3.
\]
Thus both sides of \eqref{eq:valuation} are $0$. If $r=s$, then
\eqref{eq:rec} gives $(a_u,b_u)=(a_v,b_v)$, while \eqref{eq:prefix}
gives
\[
 X_n-X_m=3(Y_v-Y_u),\qquad Y_n-Y_m=X_v-X_u.
\]
Hence the quadratic expression is multiplied by $-3$, and $n-m$ by
$3$. Repeating proves the lemma.

Let
\[
 E=\{n\ge0:(a_n,b_n)=(1,1)\}.
\]
No four points $R_n$ with $n\in E$ are collinear. Otherwise write four
of them as $R+q_iD$ $(i=0,1,2,3)$, where the $q_i$ are distinct integers
and $D=(A,B,H)$ is a primitive integer direction with $H>0$. (The integer
points on a rational line form such a one-dimensional lattice.) The lemma
applied to any pair gives
\[
 2\vthree(q_j-q_i)+\vthree(A^2-3B^2)
 =\vthree(q_j-q_i)+\vthree(H).
\]
Thus all six differences $q_j-q_i$ have one common $3$-adic valuation,
say $e$. Dividing the $q_i-q_0$ by $3^e$ would make four integers pairwise
distinct modulo $3$, which is impossible.

It remains to turn these selected vertices into a four-direction unit
walk. For $n=9q+3u+v$, with $u,v\in\{0,1,2\}$, the digit definition gives
\begin{equation}\label{eq:9}
 (a_n,b_n)=\bigl((-1)^u a_q,(-1)^v b_q\bigr).
\end{equation}
Thus the residues $r\pmod9$ for which $9q+r\in E$ are
\[
\begin{array}{c|c}
(a_q,b_q)&r\\ \hline
(1,1)&0,2,6,8\\
(1,-1)&1,7\\
(-1,1)&3,5\\
(-1,-1)&4
\end{array}
\]
Moreover, passing from $q$ to $q+1$ flips exactly one of the two signs:
in base $3$, trailing $2$'s become $0$'s, while the first digit that is
not $2$ changes parity. Hence, if
\[
 0=N_0<N_1<N_2<\cdots
\]
are the elements of $E$, then $N_{k+1}-N_k\in\{2,4,6,8\}$. Summing
$(a_j,b_j,1)$ between two consecutive elements of $E$ in the four
possible cases gives
\begin{equation}\label{eq:steps}
\begin{array}{c|cccc}
N_{k+1}-N_k&2&4&6&8\\ \hline
R_{N_{k+1}}-R_{N_k}
 &(2,0,2)&(-2,2,4)&(0,-2,6)&(-4,0,8).
\end{array}
\end{equation}

Now define $W_0=0\in\N^4$. If $N_{k+1}-N_k=2r$, let
\[
 W_{k+1}=W_k+e_r\qquad (r=1,2,3,4),
\]
where $e_1,\ldots,e_4$ are the standard basis vectors. Let
$T:\mathbb R^4\to\mathbb R^3$ be the linear map sending $e_r$ to the
$r$th displacement in \eqref{eq:steps}. By telescoping,
\[
 T(W_k)=R_{N_k}\qquad(k\ge0).
\]
If four vertices $W_{k_0},W_{k_1},W_{k_2},W_{k_3}$ were collinear,
their four images under $T$ would also be collinear. These images are
distinct because their third coordinates are
$N_{k_0}<N_{k_1}<N_{k_2}<N_{k_3}$. This contradicts the preceding
paragraph. Hence the unit-step walk $(W_k)$ has no four collinear
vertices.
\end{proof}

\end{document}
